\subsection{李多项式与代入}

设 I 为集合。设 \(\mathrm{T}_i\) 表示 I 的元素 \(i\) 在 \(\mathbf{L}(\mathbf{I})\) 中的典范像（后者有时也记作 \(\mathrm{L}((\mathrm{T}_i)_{i\in \mathbf{I}})\)）；\(\mathrm{L}(\mathbf{I})\) 的元素称为不定元 \((\mathrm{T}_i)_{i\in \mathbf{I}}\) 的李多项式。

设 g 为李代数。若 \(\boldsymbol{t} = (t_{i})_{i \in I}\) 是 g 的元素族，用 \(f_{\boldsymbol{t}}\) 表示 L(I) 到 g 的同态，使得 \(f_{\boldsymbol{t}}(\mathrm{T}_{\boldsymbol{i}}) = t_{i}\)（对 \(i \in I\)）（第 2 号，命题 1）。L(I) 的元素 P 在 \(f_{\boldsymbol{t}}\) 下的像记作 \(\mathrm{P}((t_{i})_{i \in I})\)。特别地，\(\mathrm{P}((\mathrm{T}_{\boldsymbol{i}})_{i \in I}) = \mathrm{P}\)；上述元素 \(\mathrm{P}((t_{i})_{i \in I})\) 有时称为以 \(t_{i}\) 代入 \(T_{i}\) 到李多项式 \(\mathrm{P}((\mathrm{T}_{\boldsymbol{i}})_{i \in I})\) 中所得的 g 的元素。

设 \(\sigma: \mathfrak{g} \to \mathfrak{g}'\) 为李代数同态。对每个族 \(t = (t_i)_{i \in I}\)（\(\mathfrak{g}\) 的元素的族）以及所有 \(P \in L(I)\)，

\[
\sigma (\mathrm{P} ((t _ {i}) _ {i \in \mathrm{I}})) = \mathrm{P} ((\sigma (t _ {i})) _ {i \in \mathrm{I}}),\tag{4}
\]

因为 \(\sigma^{\circ}f_{t}\) 把 \(\mathrm{T}_i\) 映为 \(\sigma(t_i)\)（对 \(i\in I\)）。

设 \((\mathbf{Q}_j)_{j\in J}\) 为 \(\mathbf{L}(\mathbf{I})\) 的元素族，并设 \(\mathbf{P}\in \mathbf{L}(\mathbf{J})\)。以 \(Q_{j}\) 代入 \(T_{j}\)（在 \(P\) 中），得到李多项式 \(R = P((Q_{j})_{j \in J}) \in L(I)\)。那么

\[
\mathrm{R} ((t _ {i}) _ {i \in \mathrm{I}}) = \mathrm{P} ((Q _ {j} ((t _ {i}) _ {i \in \mathrm{I}})) _ {j \in J})\tag{5}
\]

对每个族 \(\pmb{t} = (t_i)_{i\in \mathbf{I}}\)（李代数 \(\mathfrak{g}\) 的元素的族）成立，这可由用同态 \(f_{t}\) 作用于等式 \(\mathrm{R} = \mathrm{P}((\mathrm{Q}_j)_{j\in \mathbf{J}})\) 并利用 (4) 看出。

设 \(\mathfrak{g}\) 为李代数，I 为有限集，\(P \in L(I)\)。假设 \(\mathfrak{g}\) 是自由 K-模。映射

\[
\tilde {\mathrm{P}} \colon \mathfrak {g} ^ {\mathrm{I}} \to \mathfrak {g}
\]

由 \(\tilde{\mathrm{P}}((t_i)_{i\in \mathbf{I}}) = \mathrm{P}((t_i)_{i\in \mathbf{I}})\) 定义的映射于是是一个多项式映射。†因为 \(\mathfrak{g}^{\mathbf{I}}\) 到 \(\mathfrak{g}\) 的映射的集合 F 按由下式定义的括号成一个李代数

\[
[ \phi , \psi ] (\boldsymbol {t}) = [ \phi (\boldsymbol {t}), \psi (\boldsymbol {t}) ];\tag{6}
\]

多项式映射的集合 \(\mathbf{F}'\)（\(\mathfrak{g}^{\intercal}\) 到 \(\mathfrak{g}\) 的） 由括号的双线性是 \(\mathbf{F}\) 的李子代数。于是我们的断言由如下事实得出：映射 \(\mathrm{P} \mapsto \tilde{\mathrm{P}}\) 是李代数同态，且 \(\tilde{\mathrm{T}}_i = \mathrm{pr}_i \in \mathrm{F}'\) 对一切 \(i\) 成立。

